\section{The authority left after stopping}
An agent process can stop while its authorized work continues elsewhere. A tool
may have accepted a payment, queued a release, or begun a document update before
the operator withdraws permission. Killing the proposer does not erase those
requests. A timeout does not reveal whether they committed. Returning their
budget to a replacement agent can authorize a second effect while the first
remains possible. The relevant control question is therefore about outstanding
authority, rather than the lifetime of the model process.

This paper asks what authority can remain after stopping or withdrawal, which
mechanisms bound it, and what evidence those mechanisms support. The artifact
uses synthetic payments, releases, and edits, all within local services. These
three effect types deliberately require different observations: recipient and
amount for payments, destination and actual contents for releases, and revision
plus expected contents for edits. An authorization result concerns the encoded
permission. A task-success result concerns the intended effect. Neither result
establishes that the issuer had legitimate jurisdiction over the resource.

\subsection{Contribution and claim conditions}
We distinguish three proposed claims. First, a deterministic core makes its
specified authorization decisions independently of proposer behavior. Its
executable predicates have mechanized equivalence specifications, including
positive permissions; conservation arguments quantify over arbitrary finite
prefixes. The broker compiles this core as a Rust dependency. This verifies
particular decisions and algebraic lemmas, not the parser, database program, or
whole runtime. Their correspondence to the arguments supplied to the core is
an explicit trusted premise.

Second, revocation has two meanings. Admission revocation stops issuance of new
dispatch permits. Previously issued permits may still reach the sink and cause
effects. Their additional charge is bounded by outstanding reservations under
the cooperative sink contract. Fenced revocation adds a durable sink barrier,
which cancels outstanding old-epoch work or waits behind work that already
commits. Its acknowledgment marks a stronger boundary: subsequent old-authority
effects are excluded. We give the protocol argument, verify its decision
predicates, and exercise separate-process fault and concurrency tests.

Third, useful work under evaluated model attacks would require comparative
empirical evidence. That claim remains contingent. The release includes CPU
positive controls, descriptive controller controls, pinned public-benchmark
integration checks, and a predeclared model/inference manifest. It does not
include completed local-model campaigns. Host observations show no idle GPU on
the verified Linux hosts. No inference about attack success, checkpoint
capability, or usefulness noninferiority follows from those missing results.

\subsection{Relation to existing permission controls}
CaMeL separates control and data flow and uses explicit capabilities for agent
tool execution \cite{camel}. Progent checks symbolic tool-argument policies and
controls their evolution \cite{progent}. Permission checking by itself is
therefore not a new contribution. Fault-tolerant budget conservation already
addresses reservations, delayed effects, and authoritative settlement under
declared gateway assumptions \cite{ftbudget}. Our focus is the connection between
an executable policy core, a small separate-service implementation, and a
precise local revocation acknowledgment.

The sink contract is central to this comparison. LIMBO studies duplicate side
effects under ambiguous replies and differing service contracts \cite{limbo}.
We do not solve arbitrary external exactly-once execution. The supported sink
must atomically bind an operation identity to its normalized effect, deduplicate
it, and apply any effect in that same transaction. A deliberately non-idempotent
service is retained as an unsupported-contract negative control. Its failure is
not counted as a security failure of a contract it never implemented.

These distinctions also constrain interpretation of verification. Reference
monitor principles motivate a small mediation boundary \cite{saltzer,schneider}.
Increasing the number of checked synthetic states does not close a missing
operating-system path, validate a normalization routine, or make an incomplete
harm specification adequate. The v1 finite model and synthetic learning studies
are preserved in appendices so their evidence remains inspectable without
carrying the main revocation claim.

\subsection{Threat model}
The proposer can choose arbitrary messages, malformed requests, altered
payloads, retries, timing, and coordination among a finite number of actors.
It may try to impersonate an issuer or access databases and the sink directly.
The policy must not depend on the proposer reporting its intent, being helpful,
or recognizing revocation. Actors share one budget; copying an actor does not
copy that budget. An arbitrary finite actor population means any number of
proposers using the specified interface, not an unbounded resource-exhaustion
guarantee for a shared host.

The trusted boundary contains the session launcher, administrator channels,
normalization, the broker and sink storage adapters, the compiler and verifier,
and the Linux kernel. Administrators provision exact approvals, withdraw
issuers, and settle authoritative sink outcomes. The proposer receives none of
those commands. Storage and service processes may crash; they must preserve
committed state on recovery. We test process death, not corrupted disks,
Byzantine databases, compromised administrators, kernel exploits, or malicious
sinks. A sink that ignores deduplication or a fence invalidates the associated
bound rather than becoming an adversarial case covered by it.

We treat scheduler and storage delays as availability costs. The system does not
promise a hard deadline for revocation acknowledgment. If the sink cannot return
authoritative completion or cancellation, the broker retains the charge and the
strong acknowledgment remains unavailable. This conservative outcome is part of
the protocol, not an experimental failure to be hidden.
